\chapter{Metodologia}

Neste capítulo são apresentadas as especificações de requisitos levantadas diretamente a partir do código-fonte, a modelagem UML do sistema e a estrutura arquitetural completa do backend e banco de dados.

% ─────────────────────────────────────────────────────────────────
\section{Requisitos Funcionais}
% ─────────────────────────────────────────────────────────────────

Os requisitos funcionais foram extraídos da análise dos \textit{Controllers}, \textit{Services} e entidades do projeto:

\begin{itemize}
    \item \textbf{[RF001] Criar Conta:} O organizador deve poder se registrar na plataforma informando nome, e-mail e senha.
    \item \textbf{[RF002] Autenticar no Sistema:} O sistema deve validar as credenciais do organizador, comparando o hash BCrypt da senha armazenada, e retornar um token JWT.
    \item \textbf{[RF003] Listar e Filtrar Eventos:} O sistema deve listar apenas os eventos pertencentes ao organizador autenticado, extraindo o identificador do usuário diretamente do token JWT.
    \item \textbf{[RF004] Criar Evento:} O sistema deve permitir a criação de eventos com nome, tipo, data, horários de início e fim, capacidade máxima (\textit{Pax}) e foto de banner. A data de início não pode ser anterior ao momento da criação.
    \item \textbf{[RF005] Excluir Evento:} O sistema deve excluir logicamente (\textit{Soft Delete}) um evento, preenchendo o campo \texttt{DeletedDate} sem apagar o registro do banco.
    \item \textbf{[RF006] Gerenciar Campos Dinâmicos (EAV):} O organizador deve poder criar campos personalizados de formulário (Texto, Número ou Data), definir obrigatoriedade e reordenar os campos via operação de lote.
    \item \textbf{[RF007] Cadastrar Convidados Manualmente:} O organizador deve poder cadastrar convidados diretamente pelo painel, com validação de e-mail (expressão regular), tipo de convidado e campos dinâmicos obrigatórios do evento.
    \item \textbf{[RF008] Auto-cadastro Público de Convidados:} O sistema deve disponibilizar rotas públicas (\textit{AllowAnonymous}) para que o próprio participante preencha o formulário de inscrição, sendo bloqueado caso o horário de início do evento já tenha passado ou o limite de vagas (\textit{Pax}) tenha sido atingido.
    \item \textbf{[RF009] Check-in via QR Code (Painel Administrativo):} O aplicativo deve capturar o QR Code do participante via câmera, extrair o UUID do convidado e acionar o endpoint de check-in autenticado.
    \item \textbf{[RF010] Check-in Automático (Rota Pública):} O sistema deve disponibilizar uma rota pública decodificada de hash para que o convidado realize seu próprio check-in, respeitando o modo do evento (\texttt{CheckinEnabled}).
    \item \textbf{[RF011] Gerenciar Check-in em Lote:} O organizador deve poder aprovar ou reverter check-ins de múltiplos convidados simultaneamente via seleção em massa na interface.
    \item \textbf{[RF012] Upload de Template de Certificado:} O sistema deve aceitar arquivos Word (\texttt{.doc}/\texttt{.docx}) de até 10 MB, validar a presença do placeholder \texttt{\{Nome\}} e gerar uma pré-visualização em PDF.
    \item \textbf{[RF013] Enviar Convites por E-mail:} O sistema deve disparar e-mails com HTML personalizado e link do QR Code de auto-checkin para os convidados selecionados que ainda não fizeram check-in.
    \item \textbf{[RF014] Gerar e Enviar Certificados em Lote por E-mail:} O sistema deve substituir os \textit{placeholders} fixos e dinâmicos (EAV) no template Word, converter para PDF em memória via \textit{MemoryStream} e enviar os arquivos por e-mail para todos os convidados com check-in confirmado.
    \item \textbf{[RF015] Download de Certificados em ZIP:} O sistema deve gerar um arquivo compactado (\texttt{.zip}) contendo os PDFs individuais de todos os convidados presentes e disponibilizá-lo para download.
\end{itemize}

% ─────────────────────────────────────────────────────────────────
\section{Requisitos Não Funcionais}
% ─────────────────────────────────────────────────────────────────

\begin{itemize}
    \item \textbf{[RNF001] Autenticação JWT Stateless:} O servidor não armazena sessões. A validade e a identidade do usuário são verificadas a cada requisição pelo \textit{middleware} \texttt{[Authorize]}, que valida a assinatura HMAC-SHA256 do token.
    \item \textbf{[RNF002] Chaves Primárias UUID (Guid):} Todas as entidades herdam de \texttt{EntityBase}, que define o campo \texttt{Id} como \texttt{Guid}. O identificador é gerado pela camada de aplicação (\texttt{Guid.NewGuid()}) antes de qualquer persistência.
    \item \textbf{[RNF003] Soft Delete (Exclusão Lógica):} A classe \texttt{AuditableEntity} define os campos \texttt{CreatedDate} (preenchido automaticamente com \texttt{DateTime.UtcNow} no construtor) e \texttt{DeletedDate} (nulo por padrão). Exclusões apenas preenchem \texttt{DeletedDate}; o registro permanece intacto no banco.
    \item \textbf{[RNF004] Onion Architecture:} O projeto é divido em projetos isolados: \texttt{Domain} (entidades e interfaces puras, sem dependências externas), \texttt{Services} (regras de negócio), \texttt{Repository} (acesso a dados via EF Core), \texttt{Infra} (implementações de serviços externos) e \texttt{API} (camada de apresentação).
    \item \textbf{[RNF005] Armazenamento Seguro no Mobile:} O token JWT é persistido no dispositivo via \texttt{expo-secure-store}, que delega a criptografia ao \textit{Keychain} (iOS) ou \textit{KeyStore} (Android), sendo inacessível a outros aplicativos.
    \item \textbf{[RNF006] Processamento de Documentos em Memória:} A geração de PDF é realizada integralmente via \texttt{MemoryStream}, sem criação de arquivos temporários em disco, prevenindo acúmulo de lixo no servidor.
    \item \textbf{[RNF007] Fonte Tipográfica Consistente:} O aplicativo carrega a família tipográfica \textit{Poppins} (Regular, Bold, Black e ExtraBold) de forma assíncrona na inicialização via \texttt{expo-font}, garantindo consistência visual entre iOS e Android.
\end{itemize}

% ─────────────────────────────────────────────────────────────────
\section{Diagrama de Casos de Uso}
% ─────────────────────────────────────────────────────────────────

A Figura 1 ilustra as interações dos dois atores do sistema com as funcionalidades mapeadas a partir dos \textit{controllers}. O \textbf{Organizador} acessa as rotas protegidas pelo \texttt{[Authorize]}, enquanto o \textbf{Convidado} interage exclusivamente com o \texttt{FormController} (\texttt{[AllowAnonymous]}).

\begin{figure}[H]
\centering
\caption{Figura 1 - Diagrama de Casos de Uso do Sistema Certify}
\begin{verbatim}
```mermaid
flowchart TD
    Organizador(("Organizador"))
    Convidado(("Convidado"))

    UC0["Criar Conta"]
    UC1["Realizar Login"]
    UC2["Gerenciar Eventos"]
    UC3["Configurar Campos Dinâmicos"]
    UC4["Cadastrar Convidado Manualmente"]
    UC5["Gerenciar Check-in em Lote"]
    UC6["Scanner QR Code (Check-in Admin)"]
    UC7["Upload de Template de Certificado"]
    UC8["Enviar Convites por E-mail"]
    UC9["Enviar Certificados em Lote"]
    UC10["Download de Certificados ZIP"]
    UC11["Auto-cadastro Público"]
    UC12["Auto Check-in Online (Formulário)"]

    Organizador --> UC0
    Organizador --> UC1
    UC1 --> UC2
    UC1 --> UC3
    UC1 --> UC4
    UC1 --> UC5
    UC1 --> UC6
    UC1 --> UC7
    UC1 --> UC8
    UC1 --> UC9
    UC1 --> UC10

    Convidado --> UC11
    Convidado --> UC12
```
\end{verbatim}
\fonte{O Autor (2026)}
\end{figure}

% ─────────────────────────────────────────────────────────────────
\section{Descrição dos Casos de Uso}
% ─────────────────────────────────────────────────────────────────

\begin{table}[H]
\centering
\caption{Tabela 1 - CDU001 – Criar Conta}
\begin{tabular}{|p{4cm}|p{11cm}|}
\hline
\textbf{Identificador} & CDU001 – Criar Conta \\ \hline
\textbf{Objetivo} & Permitir que um novo organizador crie suas credenciais de acesso ao sistema. \\ \hline
\textbf{Ator principal} & Organizador \\ \hline
\textbf{Pré-condições} & Nenhuma. \\ \hline
\textbf{Pós-condições} & Registro inserido na tabela \texttt{UserProfile} com senha armazenada em hash BCrypt. \\ \hline
\textbf{Estrutura de Dados} & Nome, E-mail, Senha. \\ \hline
\multicolumn{2}{|c|}{\textbf{Fluxo Principal}} \\ \hline
\textbf{Ações do Usuário} & \textbf{Ações do Sistema} \\ \hline
1. Acessa a tela de Registro e preenche os campos Nome, E-mail e Senha. & 2. O app valida os campos localmente (formato do e-mail e comprimento mínimo da senha). \\ \hline
3. Submete o formulário. & 4. A API aplica o hash BCrypt na senha e persiste o objeto \texttt{UserProfile} no banco. Retorna HTTP 200 OK. \\ \hline
\end{tabular}
\end{table}

\begin{table}[H]
\centering
\caption{Tabela 2 - CDU002 – Realizar Login}
\begin{tabular}{|p{4cm}|p{11cm}|}
\hline
\textbf{Identificador} & CDU002 – Realizar Login \\ \hline
\textbf{Objetivo} & Autenticar o organizador e iniciar a sessão local no dispositivo. \\ \hline
\textbf{Ator principal} & Organizador \\ \hline
\textbf{Pré-condições} & Conta ativa cadastrada no sistema. \\ \hline
\textbf{Pós-condições} & Token JWT salvo no \texttt{expo-secure-store}; Redux atualizado com \texttt{signIn}; usuário redirecionado para a tela Home. \\ \hline
\textbf{Estrutura de Dados} & E-mail e Senha. \\ \hline
\multicolumn{2}{|c|}{\textbf{Fluxo Principal}} \\ \hline
\textbf{Ações do Usuário} & \textbf{Ações do Sistema} \\ \hline
1. Informa E-mail e Senha e toca em "Entrar". & 2. O app valida o formato do e-mail via regex e o comprimento mínimo da senha. \\ \hline
3. Aguarda resposta. & 4. A API consulta o \texttt{UserProfile}, compara o hash da senha com BCrypt e, se válido, assina um JWT com HMAC-SHA256 e retorna o token. \\ \hline
 & 5. O app salva o token no \texttt{expo-secure-store}, despacha \texttt{signIn} para a Redux Store e navega para a Home. \\ \hline
\multicolumn{2}{|c|}{\textbf{Fluxo Alternativo}} \\ \hline
\multicolumn{2}{|p{15cm}|}{2a. Se o e-mail ou senha forem inválidos, o app exibe um \texttt{CustomSnackBar} de erro sem realizar a chamada à API. 4a. Se as credenciais não corresponderem, a API retorna HTTP 400 e o app exibe o snackbar "E-mail ou senha incorretos".} \\ \hline
\end{tabular}
\end{table}

\begin{table}[H]
\centering
\caption{Tabela 3 - CDU004 – Criar Evento}
\begin{tabular}{|p{4cm}|p{11cm}|}
\hline
\textbf{Identificador} & CDU004 – Criar Evento \\ \hline
\textbf{Objetivo} & Registrar um novo evento vinculado ao organizador autenticado. \\ \hline
\textbf{Ator principal} & Organizador \\ \hline
\textbf{Pré-condições} & Usuário autenticado. Tipos de evento disponíveis (\texttt{EventType}). \\ \hline
\textbf{Pós-condições} & Evento persistido na tabela \texttt{Event} com o \texttt{UserId} extraído do JWT. \\ \hline
\textbf{Estrutura de Dados} & Nome, Tipo de Evento, Data, Horário Inicial, Horário Final, Capacidade (Pax), Foto. \\ \hline
\multicolumn{2}{|c|}{\textbf{Fluxo Principal}} \\ \hline
\textbf{Ações do Usuário} & \textbf{Ações do Sistema} \\ \hline
1. Acessa o formulário de novo evento e preenche os campos. & 2. O app utiliza \texttt{DateTimePicker} nativo para data/hora e \texttt{expo-image-picker} para a foto. \\ \hline
3. Confirma a criação. & 4. A API valida: nome não vazio, data/hora inicial não no passado, horário inicial menor que o final, Pax maior que zero. Persiste o evento e retorna o UUID gerado. \\ \hline
\multicolumn{2}{|c|}{\textbf{Fluxo Alternativo}} \\ \hline
\multicolumn{2}{|p{15cm}|}{4a. Se o \texttt{StartTime} for maior ou igual ao \texttt{EndTime}, a API lança \texttt{BusinessException}: "Horário inicial não pode ser maior que horário final.". 4b. Se a data e hora inicial já tiverem passado, lança: "A data e horário inicial já passaram."} \\ \hline
\end{tabular}
\end{table}

\begin{table}[H]
\centering
\caption{Tabela 4 - CDU007 – Cadastrar Convidado Manualmente}
\begin{tabular}{|p{4cm}|p{11cm}|}
\hline
\textbf{Identificador} & CDU007 – Cadastrar Convidado Manualmente \\ \hline
\textbf{Objetivo} & Registrar um participante em um evento pelo painel administrativo do aplicativo. \\ \hline
\textbf{Ator principal} & Organizador \\ \hline
\textbf{Pré-condições} & Evento ativo existente. \\ \hline
\textbf{Pós-condições} & Convidado inserido em \texttt{Guest}; respostas EAV em \texttt{EventFieldValue}; convite enviado por e-mail automaticamente. \\ \hline
\textbf{Estrutura de Dados} & Nome, E-mail, Tipo de Convidado, Foto (opcional), Campos Dinâmicos EAV. \\ \hline
\multicolumn{2}{|c|}{\textbf{Fluxo Principal}} \\ \hline
\textbf{Ações do Usuário} & \textbf{Ações do Sistema} \\ \hline
1. Preenche nome, e-mail, tipo e campos extras; envia. & 2. A API valida: nome e e-mail obrigatórios, e-mail por regex, tipo de convidado definido, campos EAV obrigatórios preenchidos e com tipo correto (Number/Date), limite de Pax não atingido, e-mail não duplicado no evento. \\ \hline
 & 3. Se tudo válido, insere \texttt{Guest} e os \texttt{EventFieldValue}; dispara convite por e-mail com link QR Code. \\ \hline
\multicolumn{2}{|c|}{\textbf{Fluxos Alternativos}} \\ \hline
\multicolumn{2}{|p{15cm}|}{2a. E-mail duplicado no evento: \textit{"Email já cadastrado para este evento."}. 2b. Limite de vagas atingido: \textit{"Limite de convidados atingido para este evento."}. 2c. Campo EAV obrigatório vazio: \textit{"\{NomeDoCampo\} é obrigatório."}.} \\ \hline
\end{tabular}
\end{table}

\begin{table}[H]
\centering
\caption{Tabela 5 - CDU009 – Check-in via QR Code (Painel Admin)}
\begin{tabular}{|p{4cm}|p{11cm}|}
\hline
\textbf{Identificador} & CDU009 – Check-in via QR Code \\ \hline
\textbf{Objetivo} & Registrar presença de um convidado por leitura óptica do QR Code. \\ \hline
\textbf{Ator principal} & Organizador \\ \hline
\textbf{Pré-condições} & Convidado cadastrado no evento; evento com \texttt{CheckinEnabled} compatível. \\ \hline
\textbf{Pós-condições} & Campo \texttt{CheckinDate} do \texttt{Guest} preenchido com o horário de Brasília. \\ \hline
\textbf{Estrutura de Dados} & UUID do convidado (embarcado no QR Code em texto puro nesta rota admin). \\ \hline
\multicolumn{2}{|c|}{\textbf{Fluxo Principal}} \\ \hline
\textbf{Ações do Usuário} & \textbf{Ações do Sistema} \\ \hline
1. Aponta a câmera para o QR Code do participante na aba "Scanner". & 2. O \texttt{CameraView} (expo-camera) detecta o código e valida se o dado está dentro dos limites do visor (\texttt{scannerBounds}). \\ \hline
3. Aguarda confirmação. & 4. O app localiza o convidado pela lista local (sem round-trip ao servidor); navega automaticamente para a tela de detalhes do convidado (\texttt{guest.js}) se encontrado. \\ \hline
 & 5. A API (\texttt{PUT /Guest/Checkin/\{id\}}) valida se o check-in já foi realizado e grava \texttt{CheckinDate = DateTime.UtcNow.ConvertToBrazilTime()}. \\ \hline
\multicolumn{2}{|c|}{\textbf{Fluxos Alternativos}} \\ \hline
\multicolumn{2}{|p{15cm}|}{2a. QR Code fora dos limites do visor: leitura ignorada. 2b. Convidado não pertence ao evento: snackbar "QRCode inválido para este evento". 5a. Check-in já realizado: API retorna \texttt{HTTP 409 Conflict}; app exibe feedback visual.} \\ \hline
\end{tabular}
\end{table}

\begin{table}[H]
\centering
\caption{Tabela 6 - CDU012 – Gerar e Enviar Certificados em Lote}
\begin{tabular}{|p{4cm}|p{11cm}|}
\hline
\textbf{Identificador} & CDU012 – Gerar e Enviar Certificados em Lote \\ \hline
\textbf{Objetivo} & Substituir os \textit{placeholders} do template Word com dados do banco, converter para PDF e enviar por e-mail. \\ \hline
\textbf{Ator principal} & Organizador \\ \hline
\textbf{Pré-condições} & Evento com template \texttt{.docx} configurado; ao menos um convidado com \texttt{CheckinDate} preenchido. \\ \hline
\textbf{Pós-condições} & Certificados PDF individuais despachados para os e-mails dos convidados presentes. \\ \hline
\textbf{Estrutura de Dados} & Placeholders fixos (\texttt{\{Nome\}, \{Data\}, \{Evento\}, \{Tipo Convidado\}, \{Tipo Evento\}, \{Horário Inicial\}, \{Horário Final\}}) e dinâmicos (campos EAV). \\ \hline
\multicolumn{2}{|c|}{\textbf{Fluxo Principal}} \\ \hline
\textbf{Ações do Usuário} & \textbf{Ações do Sistema} \\ \hline
1. Seleciona convidados presentes e aciona o envio via menu de contexto. & 2. A API filtra convidados com \texttt{CheckinDate != null} pertencentes ao lote selecionado. \\ \hline
 & 3. Para cada convidado: carrega o \texttt{.docx} via \texttt{DocX.Load()}; aplica \texttt{ReplaceText()} com \texttt{StringReplaceTextOptions} para campos fixos; itera \texttt{FieldsValues} para campos EAV dinâmicos; salva em \texttt{MemoryStream}. \\ \hline
 & 4. Converte o stream via \texttt{Spire.Doc} para PDF; monta o \texttt{MailMessageDTO} com o PDF como anexo e HTML do template \texttt{grateful.html}. \\ \hline
 & 5. Despacha a lista de e-mails em lote via \texttt{MailService.SendMailCheckfyAsync()}. \\ \hline
\multicolumn{2}{|c|}{\textbf{Fluxos Alternativos}} \\ \hline
\multicolumn{2}{|p{15cm}|}{Pré-condição não atendida: API lança \texttt{NotFoundException} com mensagem "Evento não possui template." ou "Evento não possui convidados com checkin.".} \\ \hline
\end{tabular}
\end{table}

% ─────────────────────────────────────────────────────────────────
\section{Diagrama de Classes}
% ─────────────────────────────────────────────────────────────────

O diagrama abaixo representa a hierarquia real de entidades da camada \texttt{Domain}, conforme implementada nos arquivos \texttt{.cs} do projeto. Todas as entidades auditáveis herdam de \texttt{AuditableEntity}, que por sua vez herda de \texttt{EntityBase} (que define o \texttt{Guid Id}).

\begin{figure}[H]
\centering
\caption{Figura 2 - Diagrama de Classes do Domínio}
\begin{verbatim}
```mermaid
classDiagram
    class EntityBase {
        +Guid Id
    }
    class AuditableEntity {
        +DateTime CreatedDate
        +DateTime? DeletedDate
    }
    class UserProfile {
        +String Email
        +String PasswordHash
        +String Name
        +String Photo
        +Guid ConcurrencyStamp
        +Guid SecurityStamp
    }
    class Event {
        +String Name
        +String Description
        +String Photo
        +Int? Pax
        +DateTime Date
        +TimeSpan StartTime
        +TimeSpan EndTime
        +Boolean? CheckinEnabled
        +Guid UserId
        +Guid EventTypeId
        +Guid? EventTemplateId
    }
    class EventType {
        +String Name
    }
    class EventTemplate {
        +String Path
        +String PreviewPath
    }
    class Guest {
        +String Name
        +String Email
        +String Photo
        +DateTime? CheckinDate
        +Guid EventId
        +Guid GuestTypeId
    }
    class GuestType {
        +String Name
    }
    class EventField {
        +Guid EventId
        +String Name
        +String Type
        +Boolean IsRequired
        +Int DisplayOrder
    }
    class EventFieldValue {
        +Guid EventFieldId
        +Guid GuestId
        +String Value
    }

    EntityBase <|-- AuditableEntity
    AuditableEntity <|-- UserProfile
    AuditableEntity <|-- Event
    AuditableEntity <|-- EventType
    AuditableEntity <|-- EventTemplate
    AuditableEntity <|-- Guest
    AuditableEntity <|-- GuestType
    AuditableEntity <|-- EventField
    EntityBase <|-- EventFieldValue

    UserProfile "1" --> "*" Event : organiza
    EventType "1" --> "*" Event : categoriza
    EventTemplate "0..1" --> "0..1" Event : template de
    Event "1" --> "*" Guest : contém
    Event "1" --> "*" EventField : define
    GuestType "1" --> "*" Guest : tipifica
    Guest "1" --> "*" EventFieldValue : responde
    EventField "1" --> "*" EventFieldValue : especifica
```
\end{verbatim}
\fonte{O Autor (2026)}
\end{figure}

% ─────────────────────────────────────────────────────────────────
\section{Diagramas de Sequências}
% ─────────────────────────────────────────────────────────────────

\begin{figure}[H]
\centering
\caption{Figura 3 - Diagrama de Sequência: Realizar Login}
\begin{verbatim}
```mermaid
sequenceDiagram
    actor Org as Organizador
    participant App as React Native (_main.js)
    participant SS as expo-secure-store
    participant Redux as Redux Store
    participant API as .NET 6 API (AccountController)
    participant DB as Banco de Dados

    Org->>App: Preenche e-mail e senha e toca "Entrar"
    App->>App: Valida regex e-mail + comprimento mínimo senha
    App->>API: POST /Account/Login {email, senha}
    API->>DB: SELECT UserProfile WHERE Email = ...
    DB-->>API: Retorna UserProfile com PasswordHash
    API->>API: BCrypt.Verify(senha, PasswordHash)
    API-->>App: HTTP 200 + JWT Token (HMAC-SHA256)
    App->>SS: SecureStore.setItemAsync(token)
    App->>Redux: dispatch(signIn(token))
    Redux-->>App: isSignedIn = true
    App->>Org: Monta pilha privada e navega para Home
```
\end{verbatim}
\fonte{O Autor (2026)}
\end{figure}

\begin{figure}[H]
\centering
\caption{Figura 4 - Diagrama de Sequência: Check-in via QR Code (Painel Admin)}
\begin{verbatim}
```mermaid
sequenceDiagram
    actor Org as Organizador
    participant Cam as CameraView (expo-camera)
    participant App as EventScanner.js
    participant API as GuestController (.NET 6)
    participant DB as Banco de Dados

    Org->>Cam: Aponta câmera para QR Code
    Cam-->>App: onBarcodeScanned({data, boundingBox})
    App->>App: Valida se código está dentro do scannerBounds
    App->>App: Busca convidado na lista local (guests.find)
    alt Convidado encontrado
        App->>App: navigate(routes.guest, {guest})
        App->>API: PUT /Guest/Checkin/{id}
        API->>DB: SELECT Guest WHERE Id = {id}
        DB-->>API: Retorna Guest
        API->>API: Verifica CheckinDate == null
        API->>DB: UPDATE Guest SET CheckinDate = UtcNow.ToBrazilTime()
        API-->>App: HTTP 200 OK
    else Convidado não encontrado
        App->>Org: CustomSnackBar "QRCode inválido para este evento"
    end
```
\end{verbatim}
\fonte{O Autor (2026)}
\end{figure}

\begin{figure}[H]
\centering
\caption{Figura 5 - Diagrama de Sequência: Envio de Certificados em Lote}
\begin{verbatim}
```mermaid
sequenceDiagram
    actor Org as Organizador
    participant App as guestsTab.js
    participant API as GuestController (.NET 6)
    participant Store as StorageService
    participant Xceed as Xceed.Words.NET (DocX)
    participant Spire as Spire.Doc
    participant Mail as MailService (SMTP)
    participant DB as Banco de Dados

    Org->>App: Seleciona presentes e aciona "Enviar certificado"
    App->>App: Exibe ConfirmAlert para confirmação
    App->>API: POST /Guest/Certificates/{eventId} [ids]
    API->>DB: SELECT Guests WHERE Id IN [ids] AND CheckinDate != NULL
    DB-->>API: Lista de Guests com FieldsValues
    API->>Store: Baixa arquivo .docx do StorageService
    loop Para cada Convidado
        API->>Xceed: DocX.Load(templatePath)
        Xceed->>Xceed: ReplaceText({Nome}, {Data}, {Tipo Convidado}...)
        Xceed->>Xceed: ReplaceText(EAV: {fieldValue.EventField.Name})
        Xceed->>API: MemoryStream do DOCX editado
        API->>Spire: ConvertDocumentToPdf(documentStream)
        Spire-->>API: pdfStream
        API->>API: Monta MailMessageDTO + AddAttachment(pdfStream)
    end
    API->>Mail: SendMailCheckfyAsync(List<MailMessageDTO>)
    Mail-->>Org: E-mails entregues com PDF anexado
    API-->>App: HTTP 200 OK
```
\end{verbatim}
\fonte{O Autor (2026)}
\end{figure}

% ─────────────────────────────────────────────────────────────────
\section{Visão Geral do Projeto do Sistema}
% ─────────────────────────────────────────────────────────────────

\begin{figure}[H]
\centering
\caption{Figura 6 - Arquitetura Distribuída: Cliente, Servidor e Banco de Dados}
\begin{verbatim}
```mermaid
graph TD
    subgraph Mobile ["App Mobile (React Native + Expo SDK 51)"]
        UI["Telas / Componentes"] --> Redux["Redux Store (token, snackBar)"]
        Redux --> SS["expo-secure-store (JWT)"]
        UI --> Axios["Axios Client\n(Header: Bearer JWT)"]
    end

    subgraph API ["Backend (.NET 6 - Onion Architecture)"]
        Controllers["API Layer (Controllers)"]
        Services["Services Layer (Regras de Negócio)"]
        Repos["Repository Layer (EF Core)"]
        DocEngine["Xceed.Words.NET + Spire.Doc"]
        MailSvc["MailService (SMTP)"]
        QRSvc["QRCodeService"]

        Controllers --> Services
        Services --> Repos
        Services --> DocEngine
        Services --> MailSvc
        Services --> QRSvc
    end

    subgraph DB ["Banco de Dados Relacional"]
        EF["Entity Framework Core"]
        SGBD[("SGBD - Chaves UUID")]
        EF --> SGBD
    end

    Axios <-->|"REST HTTPS / JSON"| Controllers
    Repos --> EF
```
\end{verbatim}
\fonte{O Autor (2026)}
\end{figure}

% ─────────────────────────────────────────────────────────────────
\section{Visão do Banco de Dados}
% ─────────────────────────────────────────────────────────────────

\subsection{SGBD}

O sistema utiliza o \textbf{Entity Framework Core} como ORM (\textit{Object-Relational Mapper}), que abstrai as consultas SQL e permite o mapeamento direto entre as entidades C\# e as tabelas relacionais. A adoção de chaves primárias \texttt{Guid} torna a aplicação agnóstica ao SGBD, sendo compatível com SQL Server, PostgreSQL e SQLite sem alteração de código no domínio.

\subsection{MER - Modelo de Entidade Relacional}

\begin{figure}[H]
\centering
\caption{Figura 7 - Modelo de Entidade Relacional (MER)}
\begin{verbatim}
```mermaid
erDiagram
    UserProfile {
        uuid Id PK
        varchar Email
        varchar PasswordHash
        varchar Name
        text Photo
        uuid ConcurrencyStamp
        uuid SecurityStamp
        timestamp CreatedDate
        timestamp DeletedDate
    }
    Event {
        uuid Id PK
        uuid UserId FK
        uuid EventTypeId FK
        uuid EventTemplateId FK
        varchar Name
        text Description
        text Photo
        int Pax
        date Date
        time StartTime
        time EndTime
        boolean CheckinEnabled
        timestamp CreatedDate
        timestamp DeletedDate
    }
    EventType {
        uuid Id PK
        varchar Name
        timestamp CreatedDate
        timestamp DeletedDate
    }
    EventTemplate {
        uuid Id PK
        varchar Path
        varchar PreviewPath
        timestamp CreatedDate
        timestamp DeletedDate
    }
    Guest {
        uuid Id PK
        uuid EventId FK
        uuid GuestTypeId FK
        varchar Name
        varchar Email
        text Photo
        timestamp CheckinDate
        timestamp CreatedDate
        timestamp DeletedDate
    }
    GuestType {
        uuid Id PK
        varchar Name
    }
    EventField {
        uuid Id PK
        uuid EventId FK
        varchar Name
        varchar Type
        boolean IsRequired
        smallint DisplayOrder
        timestamp CreatedDate
        timestamp DeletedDate
    }
    EventFieldValue {
        uuid Id PK
        uuid EventFieldId FK
        uuid GuestId FK
        text Value
    }

    UserProfile ||--o{ Event : "organiza"
    EventType ||--o{ Event : "categoriza"
    EventTemplate |o--o{ Event : "template de"
    Event ||--o{ Guest : "contém"
    Event ||--o{ EventField : "define campos"
    GuestType ||--o{ Guest : "tipifica"
    Guest ||--o{ EventFieldValue : "responde"
    EventField ||--o{ EventFieldValue : "especifica"
```
\end{verbatim}
\fonte{O Autor (2026)}
\end{figure}

% ─────────────────────────────────────────────────────────────────
\section{Visão do Backend (API RESTful)}
% ─────────────────────────────────────────────────────────────────

\subsection{Tecnologias a serem usadas}

\begin{itemize}
    \item \textbf{ASP.NET Core 6 (C\#):} Framework principal da API REST, responsável pelo roteamento, injeção de dependências e \textit{middleware} de autenticação JWT.
    \item \textbf{Entity Framework Core:} ORM responsável pelo mapeamento objeto-relacional e execução de queries via LINQ.
    \item \textbf{Xceed.Words.NET (DocX):} Biblioteca para manipulação de documentos \texttt{.docx} em memória. Usada para carregar o template e realizar substituições via \texttt{StringReplaceTextOptions}.
    \item \textbf{Spire.Doc:} Biblioteca de conversão de documentos Word para PDF, operando integralmente via \texttt{MemoryStream} sem geração de arquivos temporários.
    \item \textbf{JWT Bearer Authentication:} Middleware \texttt{[Authorize]} que valida a assinatura HMAC-SHA256 do token em todas as rotas protegidas.
\end{itemize}

\subsection{Projeto do Backend (Onion Architecture)}

A API é estruturada em cinco projetos isolados que seguem a \textit{Onion Architecture}:

\begin{itemize}
    \item \textbf{Domain:} Contém entidades (\texttt{Event}, \texttt{Guest}, etc.), interfaces de serviços e repositórios, DTOs e exceções de negócio (\texttt{BusinessException}, \texttt{NotFoundException}, \texttt{ConflictException}). Sem nenhuma dependência externa.
    \item \textbf{Services:} Implementações das regras de negócio (validações, geração de QR Code, substituição de \textit{placeholders}, envio de e-mail). Dependem apenas de interfaces do \texttt{Domain}.
    \item \textbf{Repository:} Implementações dos repositórios com \texttt{EF Core}. Aplicam o filtro de \textit{Soft Delete} nas consultas.
    \item \textbf{Infra:} Implementações de serviços externos (armazenamento de arquivos, e-mail SMTP, hash de senhas).
    \item \textbf{API:} Camada de apresentação. Contém os \textit{Controllers} e o \textit{pipeline} de configuração do ASP.NET Core.
\end{itemize}

\subsection{Endpoints da API}

\begin{table}[H]
\centering
\caption{Tabela 7 - Mapeamento Completo de Endpoints da API Certify}
\footnotesize
\begin{tabular}{|l|l|p{5cm}|l|}
\hline
\textbf{Método} & \textbf{Rota} & \textbf{Descrição} & \textbf{Auth} \\ \hline
POST & /Account/Login & Autentica e retorna JWT & Não \\ \hline
GET & /Event/GetEvents & Lista eventos do usuário logado & Sim \\ \hline
GET & /Event/GetEvent/\{id\} & Retorna evento com Guests e Fields & Sim \\ \hline
POST & /Event/NewEvent & Cria novo evento & Sim \\ \hline
DELETE & /Event/Delete/\{id\} & Soft delete do evento & Sim \\ \hline
PUT & /Event/CheckinEnabledMode & Altera o modo de check-in & Sim \\ \hline
GET & /Event/QRCode/\{eventId\} & Gera QR Code do link de auto check-in & Não \\ \hline
GET & /Event/QRCode/decode/\{eventId\} & Decodifica hash do evento & Não \\ \hline
GET & /Event/Certificado/Download/\{id\} & Gera ZIP com todos os PDFs & Sim \\ \hline
POST & /Event/Certificado/Send/\{id\} & Envia certificados por e-mail & Sim \\ \hline
GET & /Event/\{id\}/EventField & Lista campos EAV do evento & Sim \\ \hline
POST & /Event/EventField & Cria novo campo dinâmico & Sim \\ \hline
PUT & /Event/EventField/Reorder & Reordena campos (DTO) & Sim \\ \hline
PUT & /Event/\{id\}/EventField/Reorder & Reordena campos (array Guid) & Sim \\ \hline
DELETE & /EventField/\{id\} & Remove campo dinâmico & Sim \\ \hline
POST & /EventTemplate/Upload/\{id\} & Upload do template .docx & Sim \\ \hline
DELETE & /EventTemplate/Remove/\{id\} & Remove template do evento & Sim \\ \hline
GET & /EventType/GetEventTypes & Lista tipos de eventos & Sim \\ \hline
GET & /GuestType/GetGuestTypes & Lista tipos de convidados & Sim \\ \hline
POST & /Guest/NewGuest & Cadastra convidado (admin) & Sim \\ \hline
PUT & /Guest/Checkin/\{id\} & Check-in individual & Sim \\ \hline
PUT & /Guest/Uncheckin/\{id\} & Reverte check-in individual & Sim \\ \hline
PUT & /Guest/Checkin & Check-in em lote [array Guid] & Sim \\ \hline
PUT & /Guest/Uncheckin & Reverte check-in em lote & Sim \\ \hline
DELETE & /Guest/Delete/\{id\} & Soft delete de convidado & Sim \\ \hline
DELETE & /Guest/Delete & Soft delete em lote & Sim \\ \hline
POST & /Guest/Invitations/\{eventId\} & Envia convites por e-mail em lote & Sim \\ \hline
POST & /Guest/Certificates/\{eventId\} & Gera e envia certificados selecionados & Sim \\ \hline
GET & /Guest/QRCode/\{guestId\} & Gera QR Code do convidado (hash) & Não \\ \hline
GET & /Form/Guest/\{eventId\} & Formulário público de inscrição & Não \\ \hline
POST & /Form/NewGuest & Submete inscrição pública & Não \\ \hline
GET & /Form/Checkin/\{eventId\} & Formulário público de check-in & Não \\ \hline
GET & /Form/Guest/Search/\{code\} & Busca convidado por código de acesso & Não \\ \hline
PUT & /Form/Guest/Checkin/\{code\} & Auto check-in via código de acesso & Não \\ \hline
\end{tabular}
\end{table}
